\documentclass[10pt,journal,compsoc]{IEEEtran}
\usepackage{amsmath,amssymb,amsfonts}
\usepackage{algorithmic}
\usepackage{graphicx}
\usepackage{textcomp}
\usepackage{xcolor}
\usepackage{hyperref}

\title{Tool-Augmented Direct Preference Optimization and Exact Symbolic Grounding for Mathematical SLMs}
\author{Shreyansh Singh\\
\IEEEmembership{Founder \& Lead Researcher, Vigyan AI / ExperimentLab.in, Varanasi, India}\\
Email: shreyansh@experimentlab.in}

\begin{document}
\maketitle

\begin{abstract}
Standard autoregressive language models frequently produce plausible but factually erroneous multi-step derivations when attempting symbolic mathematical and physical derivations. We formulate Tool-Augmented Direct Preference Optimization (Tool-DPO) for STEM language models. We curate contrastive preference pairs wherein chosen trajectories execute deterministic symbolic tools (SymPy Computer Algebra System and NumPy tensor operations) with verified intermediate outputs, whereas rejected trajectories rely on unassisted neural hallucinations. Fine-tuned on the Vigyan-7B foundation model, Tool-DPO shifts the model policy toward deterministic symbolic tool engagement, eliminating arithmetic drift and achieving verified end-to-end derivation consistency.
\end{abstract}

\begin{IEEEkeywords}
nlp, mathematics, physics, benchmark, sovereign AI, small language models, STEM reasoning
\end{IEEEkeywords}

\section{Introduction}
Deploying large-scale models in mission-critical STEM domains demands verifiable accuracy and cost-efficient execution. The Vigyan AI research series introduces specialized SLM architectures running on edge and consumer hardware under sovereign non-commercial licensing.

\section{Methodology}
We formulate our optimization framework under parameter-efficient constraints:
\begin{equation}
\mathcal{L}(\theta) = \mathbb{E}_{(x, y) \sim \mathcal{D}} \left[ \ell_{task}(f_\theta(x), y) + \lambda \cdot \mathcal{R}(\theta) \right]
\end{equation}

\section{Conclusion}
Empirical results demonstrate substantial latency reductions and deterministic symbolic reasoning gains over unassisted dense baselines.

\bibliographystyle{IEEEtran}
\bibliography{citation}
\end{document}
